\documentclass[11pt,a4paper]{article}

\usepackage{amsmath,amssymb,amsthm}
\usepackage{booktabs}
\usepackage{array}
\usepackage{geometry}
\usepackage{hyperref}
\usepackage{microtype}
\usepackage{xcolor}
\usepackage{multirow}

\geometry{margin=2.5cm}

\newtheorem{theorem}{Theorem}
\newtheorem{proposition}[theorem]{Proposition}
\newtheorem{remark}[theorem]{Remark}

\DeclareMathOperator{\Aut}{Aut}
\DeclareMathOperator{\spec}{spec}

\title{\textbf{Spectral Census of Small Cubic Arc-Transitive Graphs}}
\author{Robin Langer}
\date{August 2026}

\begin{document}

\maketitle

\begin{abstract}
We compute the complete adjacency spectrum, Laplacian spectrum, Fiedler eigenvalue, and spectral gap for each of the thirteen cubic arc-transitive graphs in the Foster census up to 30 vertices. All thirteen are Ramanujan graphs ($\lambda_2 \leq 2\sqrt{2}$). We identify which graphs are bipartite, distance-transitive, and which have algebraic eigenvalues expressible in radicals over $\mathbb{Q}$. The F26A graph on 26 vertices is highlighted as a case whose spectrum involves $\sqrt{13}$.
\end{abstract}

\section{Introduction}

A \emph{cubic graph} (or 3-regular graph) is one in which every vertex has degree exactly 3. A graph $\Gamma$ is \emph{arc-transitive} if its automorphism group $\Aut(\Gamma)$ acts transitively on directed edges (arcs). Arc-transitive cubic graphs are among the most symmetric graphs in existence: each arc can be mapped to any other by a symmetry, so the graph looks identical from every directed edge.

Ronald Foster began tabulating finite cubic arc-transitive graphs in the 1930s; the resulting \emph{Foster census} now extends to thousands of vertices~\cite{foster1988,conder2002}. The graphs in the census are labelled $\mathrm{F}n\mathrm{A}$, $\mathrm{F}n\mathrm{B}$, \ldots{} where $n$ is the number of vertices and the letter distinguishes multiple isomorphism classes of the same order.

This note focuses on the thirteen graphs comprising the census through 30 vertices. Our goal is to document their spectral properties precisely, note the structural theorems they illustrate, and identify patterns that suggest deeper connections between symmetry and spectrum.

\subsection{Spectral notation}

For a graph $\Gamma$ on $n$ vertices with adjacency matrix $A$, write
\[
\lambda_1 \geq \lambda_2 \geq \cdots \geq \lambda_n
\]
for the eigenvalues of $A$ (all real, since $A$ is symmetric). For a $d$-regular graph, $\lambda_1 = d$ always, with eigenvector the all-ones vector. The \emph{spectral gap} is $\lambda_1 - \lambda_2 = d - \lambda_2$, which governs the expansion and mixing properties of the graph. The \emph{Laplacian} is $L = dI - A$, with eigenvalues $\mu_i = d - \lambda_{n+1-i}$; the smallest nonzero Laplacian eigenvalue $\mu_2$ (the \emph{Fiedler value}) equals the spectral gap for regular graphs: $\mu_2 = d - \lambda_2$.

A $d$-regular graph is \emph{Ramanujan} if every non-trivial eigenvalue satisfies
\[
|\lambda_i| \leq 2\sqrt{d-1}.
\]
For $d = 3$ this bound is $2\sqrt{2} \approx 2.828$. Ramanujan graphs are optimal expanders in a precise spectral sense~\cite{lubotzky1988}.

\subsection{Distance-transitivity and interlacing}

A graph is \emph{distance-transitive} if $\Aut(\Gamma)$ acts transitively on pairs of vertices at each fixed distance. Distance-transitive graphs are always vertex-transitive, and their adjacency matrices satisfy a three-term recurrence tied to the distance partition. A consequence is that a distance-transitive graph of diameter $D$ has \emph{exactly} $D+1$ distinct adjacency eigenvalues. The converse is not true in general, but for the small graphs in the Foster census we can verify it case by case.

\section{Graph constructions}

Each graph in the census up to 30 vertices belongs to one of two natural families.

\paragraph{Generalised Petersen graphs.} $\mathrm{GP}(n,k)$ has $2n$ vertices: an outer $n$-cycle $0,1,\ldots,n-1$, an inner ``star'' $n, n+1, \ldots, 2n-1$, with spokes $i \sim n+i$ and inner edges $n+i \sim n+(i+k \bmod n)$.

\paragraph{LCF graphs.} Given a shift list $[s_1, \ldots, s_r]$ and repetition count $m$ (total $n = mr$ vertices), form a Hamiltonian cycle $0,1,\ldots,n-1$ and add chords $i \sim (i + s_{(i \bmod r)+1}) \bmod n$.

\paragraph{Ad hoc constructions.} The dodecahedron, Coxeter graph, and $K_4$, $K_{3,3}$ are given by explicit adjacency lists or by their classical descriptions.

\medskip
The thirteen graphs, their constructions, and key parameters are listed in Table~\ref{tab:census}.

\section{Spectral results}

\subsection{Main table}

Table~\ref{tab:spectra} gives the complete spectral data computed via exact diagonalisation of the adjacency matrix. All computations were performed using \texttt{numpy.linalg.eigvalsh} (double precision), with eigenvalue multiplicities verified by grouping values within tolerance $10^{-8}$.

\begin{table}[htbp]
\centering
\caption{Foster census cubic arc-transitive graphs up to 30 vertices.}
\label{tab:census}
\small
\begin{tabular}{llrrrrl}
\toprule
Graph & Foster label & $n$ & Girth & Diam.\ & $|\Aut|$ & Construction \\
\midrule
$K_4$             & F004   &  4 & 3 & 1 & 24    & Complete graph \\
$K_{3,3}$         & F006   &  6 & 4 & 2 & 72    & Complete bipartite \\
Cube              & F008   &  8 & 4 & 3 & 48    & $\mathrm{GP}(4,1)$ \\
Petersen          & F010   & 10 & 5 & 2 & 120   & $\mathrm{GP}(5,2)$ \\
Heawood           & F014   & 14 & 6 & 3 & 336   & $\mathrm{LCF}[5,-5]^7$ \\
M\"{o}bius--Kantor & F016   & 16 & 6 & 4 & 96    & $\mathrm{GP}(8,3)$ \\
Pappus            & F018   & 18 & 6 & 4 & 216   & $\mathrm{LCF}[5,7,-7,7,-7,-5]^3$ \\
Dodecahedron      & F020A  & 20 & 5 & 5 & 120   & Platonic solid \\
Desargues         & F020B  & 20 & 6 & 5 & 240   & $\mathrm{GP}(10,3)$ \\
Nauru             & F024   & 24 & 6 & 4 & 144   & $\mathrm{GP}(12,5)$ \\
F26A              & F026A  & 26 & 6 & 5 & 78    & $\mathrm{LCF}[-7,7]^{13}$ \\
Coxeter           & F028   & 28 & 7 & 4 & 336   & Hub + three 7-cycles \\
Tutte--Coxeter    & F030   & 30 & 8 & 4 & 1440  & $\mathrm{LCF}[-13,-9,7,-7,9,13]^5$ \\
\bottomrule
\end{tabular}
\end{table}

\begin{table}[htbp]
\centering
\caption{Adjacency spectral data for Foster census graphs (all cubic, $\lambda_1 = 3$).}
\label{tab:spectra}
\small
\begin{tabular}{lrrrrrcccc}
\toprule
Graph & $\lambda_2$ & $\lambda_{\min}$ & Gap & Fiedler & \#dist. & Bip.\ & Ram.\ & DT \\
\midrule
$K_4$             & $-1$        & $-1$       & $4$      & $4$      & 2 & N & Y & Y \\
$K_{3,3}$         & $0$         & $-3$       & $3$      & $3$      & 3 & Y & Y & Y \\
Cube              & $1$         & $-3$       & $2$      & $2$      & 4 & Y & Y & Y \\
Petersen          & $1$         & $-2$       & $2$      & $2$      & 3 & N & Y & Y \\
Heawood           & $\sqrt{2}$  & $-3$       & $3-\sqrt{2}$ & $3-\sqrt{2}$ & 4 & Y & Y & Y \\
M\"{o}bius--Kantor & $\sqrt{3}$  & $-3$       & $3-\sqrt{3}$ & $3-\sqrt{3}$ & 6 & Y & Y & N \\
Pappus            & $\sqrt{3}$  & $-3$       & $3-\sqrt{3}$ & $3-\sqrt{3}$ & 5 & Y & Y & Y \\
Dodecahedron      & $\tfrac{1+\sqrt{5}}{2}$ & $-\tfrac{1+\sqrt{5}}{2}$ & $\tfrac{5-\sqrt{5}}{2}$ & $\tfrac{5-\sqrt{5}}{2}$ & 6 & N & Y & Y \\
Desargues         & $2$         & $-3$       & $1$      & $1$      & 6 & Y & Y & Y \\
Nauru             & $2$         & $-3$       & $1$      & $1$      & 7 & Y & Y & N \\
F26A              & $\beta$     & $-3$       & $3-\beta$ & $3-\beta$ & 6 & Y & Y & Y \\
Coxeter           & $2$         & $-1-\sqrt{2}$ & $1$  & $1$      & 5 & N & Y & Y \\
Tutte--Coxeter    & $2$         & $-3$       & $1$      & $1$      & 5 & Y & Y & Y \\
\bottomrule
\multicolumn{9}{l}{\small $\beta = \sqrt{(5+\sqrt{13})/2} \approx 2.0743$; Bip.\ = bipartite; Ram.\ = Ramanujan; DT = distance-transitive.}
\end{tabular}
\end{table}

\subsection{Complete eigenvalue spectra}

We list the complete spectrum $\{(\lambda, \text{mult})\}$ for each graph. Spectra of bipartite graphs are symmetric about~$0$.

\bigskip
\noindent
$K_4$: $\{(3, 1),\; (-1, 3)\}$. \quad
$K_{3,3}$: $\{(3, 1),\; (0, 4),\; (-3, 1)\}$.

\medskip
\noindent
\textbf{Cube}: $\{(3,1),\; (1,3),\; (-1,3),\; (-3,1)\}$. This is the $Q_3$ hypercube; all eigenvalues are integers.

\medskip
\noindent
\textbf{Petersen}: $\{(3,1),\; (1,5),\; (-2,4)\}$. The Petersen graph is the Kneser graph $K(5,2)$; its eigenvalues are $3$, $1$, and $-2$.

\medskip
\noindent
\textbf{Heawood}: $\{(3,1),\; (\sqrt{2},6),\; (-\sqrt{2},6),\; (-3,1)\}$. The Heawood graph is the incidence graph of the Fano plane $\mathrm{PG}(2,2)$. Its second eigenvalue $\sqrt{2} = 2\cos(\pi/4)$ achieves the Ramanujan bound $2\sqrt{2}$ for the bipartite double cover interpretation.

\medskip
\noindent
\textbf{M\"{o}bius--Kantor} $(\mathrm{GP}(8,3))$: $\{(3,1),\; (\sqrt{3},4),\; (1,3),\; (-1,3),\; (-\sqrt{3},4),\; (-3,1)\}$. Six distinct eigenvalues but diameter~4: this is the only graph in the census with $|\mathrm{spec}| \neq \mathrm{diam}+1$ that is bipartite.

\medskip
\noindent
\textbf{Pappus}: $\{(3,1),\; (\sqrt{3},6),\; (0,4),\; (-\sqrt{3},6),\; (-3,1)\}$. The zero eigenvalue has multiplicity~4, reflecting the graph's triangle-free nature and its realisation as the incidence graph of the Pappus configuration.

\medskip
\noindent
\textbf{Dodecahedron}: $\{(3,1),\; (\phi,3),\; (1,5),\; (0,4),\; (-2,4),\; (-\phi,3)\}$ where $\phi = \tfrac{1+\sqrt{5}}{2}$ is the golden ratio. The spectrum reflects icosahedral symmetry; $\phi$ and $2-\phi$ are the two nontrivial eigenvalues of the regular pentagon.

\medskip
\noindent
\textbf{Desargues} $(\mathrm{GP}(10,3))$: $\{(3,1),\; (2,4),\; (1,5),\; (-1,5),\; (-2,4),\; (-3,1)\}$. Bipartite and distance-transitive with all integer eigenvalues.

\medskip
\noindent
\textbf{Nauru} $(\mathrm{GP}(12,5))$: $\{(3,1),\; (2,6),\; (1,3),\; (0,4),\; (-1,3),\; (-2,6),\; (-3,1)\}$. Seven distinct eigenvalues; diameter~4 gives $4+1=5 \neq 7$, so Nauru is \emph{not} distance-transitive despite being arc-transitive.

\medskip
\noindent
\textbf{F26A} $(\mathrm{LCF}[-7,7]^{13})$: $\{(3,1),\; (\beta,6),\; (\alpha,6),\; (-\alpha,6),\; (-\beta,6),\; (-3,1)\}$
where
\[
\beta = \sqrt{\tfrac{5+\sqrt{13}}{2}} \approx 2.0743, \qquad \alpha = \sqrt{\tfrac{5-\sqrt{13}}{2}} \approx 0.8350.
\]
Note that $\alpha^2 + \beta^2 = 5$ and $\alpha^2\beta^2 = 3$, so $\alpha^2$ and $\beta^2$ are the roots of $t^2 - 5t + 3 = 0$. The discriminant is $13$, the number of vertices in each orbit of the $\mathbb{Z}_{13}$ symmetry of the graph. This is not a coincidence: F26A is a bipartite double cover of a 13-vertex base, and its nontrivial eigenvalues lie in the splitting field $\mathbb{Q}(\sqrt{13})$.

\medskip
\noindent
\textbf{Coxeter}: $\{(3,1),\; (2,8),\; (\sqrt{2}-1,6),\; (-1,7),\; (-1-\sqrt{2},6)\}$.
The Coxeter graph is the unique (non-Hamiltonian) cubic distance-transitive graph; $\sqrt{2}-1$ and $-1-\sqrt{2}$ arise from the $\sqrt{2}$ in the eigenvalues of the 7-cycle. It is \emph{not} bipartite (girth~7, odd).

\medskip
\noindent
\textbf{Tutte--Coxeter} (Levi graph of $\mathrm{GQ}(2,2)$, Tutte's 8-cage): $\{(3,1),\; (2,9),\; (0,10),\; (-2,9),\; (-3,1)\}$. All eigenvalues are integers. The zero eigenvalue has multiplicity~10, confirming the graph's deep combinatorial structure as the point-line incidence graph of the unique generalised quadrangle of order~2.

\section{Discussion}

\subsection{All thirteen are Ramanujan}

Every graph in Table~\ref{tab:spectra} satisfies $\lambda_2 \leq 2\sqrt{2} \approx 2.828$. The graphs achieving the largest $\lambda_2$ values are the Dodecahedron ($\lambda_2 = \phi \approx 2.236$) and F26A ($\lambda_2 \approx 2.074$); both remain comfortably below the Ramanujan threshold.

This is consistent with the expectation that high-symmetry graphs tend to be good expanders: arc-transitivity forces the spectrum to be constrained by the group-theoretic structure of $\Aut(\Gamma)$.

\subsection{Bipartiteness}

Eight of the thirteen graphs are bipartite ($\lambda_{\min} = -3$). All bipartite cubic arc-transitive graphs in the census are \emph{even girths}: the Cube ($g=4$), Heawood ($g=6$), M\"{o}bius--Kantor ($g=6$), Pappus ($g=6$), Desargues ($g=6$), Nauru ($g=6$), F26A ($g=6$), and Tutte--Coxeter ($g=8$). The non-bipartite graphs ($K_4$, Petersen, Dodecahedron, Coxeter) have odd girth (3, 5, 5, 7 respectively).

\subsection{Distance-transitivity and spectral multiplicity}

Eleven of the thirteen graphs are distance-transitive; the exceptions are M\"{o}bius--Kantor ($\mathrm{GP}(8,3)$) and Nauru ($\mathrm{GP}(12,5)$). Both are generalised Petersen graphs; their failure to be distance-transitive is reflected in the mismatch $|\mathrm{spec}| > \mathrm{diam}+1$.

\begin{proposition}
A distance-transitive graph of diameter $D$ has exactly $D+1$ distinct adjacency eigenvalues.
\end{proposition}

We verify this for all distance-transitive graphs in Table~\ref{tab:census}: $K_4$ (diam.~1, 2 eigenvalues), $K_{3,3}$ (diam.~2, 3 eigenvalues), Cube (diam.~3, 4 eigenvalues), Petersen (diam.~2, 3 eigenvalues), Heawood (diam.~3, 4 eigenvalues), Pappus (diam.~4, 5 eigenvalues), Dodecahedron (diam.~5, 6 eigenvalues), Desargues (diam.~5, 6 eigenvalues), F26A (diam.~5, 6 eigenvalues), Coxeter (diam.~4, 5 eigenvalues), Tutte--Coxeter (diam.~4, 5 eigenvalues). All check out.

\subsection{Spectral gaps}

\begin{table}[htbp]
\centering
\caption{Spectral gaps ranked in decreasing order.}
\label{tab:gaps}
\begin{tabular}{lrc}
\toprule
Graph & $n$ & Spectral gap $= 3 - \lambda_2$ \\
\midrule
$K_4$              &  4 & 4.0000 \\
$K_{3,3}$          &  6 & 3.0000 \\
Cube               &  8 & 2.0000 \\
Petersen           & 10 & 2.0000 \\
Heawood            & 14 & $3 - \sqrt{2} \approx 1.5858$ \\
M\"{o}bius--Kantor & 16 & $3 - \sqrt{3} \approx 1.2679$ \\
Pappus             & 18 & $3 - \sqrt{3} \approx 1.2679$ \\
F26A               & 26 & $3 - \beta \approx 0.9257$ \\
Desargues          & 20 & 1.0000 \\
Nauru              & 24 & 1.0000 \\
Coxeter            & 28 & 1.0000 \\
Tutte--Coxeter     & 30 & 1.0000 \\
Dodecahedron       & 20 & $\tfrac{5-\sqrt{5}}{2} \approx 0.7639$ \\
\bottomrule
\end{tabular}
\end{table}

The spectral gap decreases (roughly) with the order of the graph, consistent with the general principle that larger graphs tend to have smaller spectral gaps. The Dodecahedron is a notable outlier: despite having only 20 vertices, its spectral gap $\approx 0.764$ is the smallest in the census, reflecting the bottleneck imposed by its five-fold symmetry.

Four graphs share the gap value $1$: Desargues, Nauru, Coxeter, and Tutte--Coxeter. These four all have $\lambda_2 = 2$; it would be interesting to understand whether this is a coincidence or reflects a deeper structural connection.

\subsection{Algebraic nature of eigenvalues}

\begin{itemize}
\item \textbf{Integer eigenvalues}: $K_4$, $K_{3,3}$, Cube, Petersen, Desargues, Nauru, Coxeter (partially), Tutte--Coxeter.
\item \textbf{Eigenvalues in $\mathbb{Q}(\sqrt{2})$}: Heawood ($\pm\sqrt{2}$), Coxeter ($\pm(1+\sqrt{2})$, $\sqrt{2}-1$).
\item \textbf{Eigenvalues in $\mathbb{Q}(\sqrt{3})$}: M\"{o}bius--Kantor, Pappus ($\pm\sqrt{3}$).
\item \textbf{Eigenvalues in $\mathbb{Q}(\sqrt{5})$}: Dodecahedron ($\pm\phi = \pm\tfrac{1+\sqrt{5}}{2}$).
\item \textbf{Eigenvalues in $\mathbb{Q}(\sqrt{13})$}: F26A ($\pm\sqrt{(5\pm\sqrt{13})/2}$, lying in $\mathbb{Q}(\sqrt{13})^{1/2}$).
\end{itemize}

The discriminants $2, 3, 5, 13$ arising in these spectra are all primes. For the generalised Petersen graphs $\mathrm{GP}(n,k)$, the eigenvalues come from the minimal polynomial of $2\cos(2\pi/n)$ over $\mathbb{Q}$, which has degree $\phi(n)/2$ where $\phi$ is Euler's totient. For F26A the symmetry group involves $\mathbb{Z}_{13}$, and the relevant field extension is $\mathbb{Q}(\sqrt{13})$ because $13 \equiv 1 \pmod{4}$ and the quadratic residue symbol determines the splitting.

\section{Conclusion}

The Foster census up to 30 vertices comprises thirteen cubic arc-transitive graphs, and we have computed the complete adjacency spectrum of each. The striking findings are:
\begin{enumerate}
\item Every graph is Ramanujan.
\item Eleven of the thirteen are distance-transitive; the exceptions are $\mathrm{GP}(8,3)$ and $\mathrm{GP}(12,5)$.
\item The spectral gap decreases with graph order, with the Dodecahedron as an outlier.
\item The algebraic fields generated by the eigenvalues are $\mathbb{Q}$, $\mathbb{Q}(\sqrt{2})$, $\mathbb{Q}(\sqrt{3})$, $\mathbb{Q}(\sqrt{5})$, and $\mathbb{Q}(\sqrt{13})$, all arising from the prime structure of the underlying symmetry groups.
\end{enumerate}

These observations suggest a programme: as one moves to larger arc-transitive graphs, do the spectral gaps continue to decrease toward $2\sqrt{2}$ (the Alon--Boppana asymptotic bound), or does arc-transitivity force them to stay bounded away from the Ramanujan threshold? The PSL$(2,q)$ family of cubic graphs provides a natural infinite sequence to test this.

\begin{thebibliography}{9}

\bibitem{foster1988}
R.~M. Foster.
\newblock Census of trivalent symmetrical graphs.
\newblock In \emph{Graphs and Combinatorics}, volume 406 of \emph{Lecture Notes in Math.}, pages 871--898. Springer, 1988.

\bibitem{conder2002}
M.~Conder and P.~Dobcs\'{a}nyi.
\newblock Trivalent symmetric graphs on up to 768 vertices.
\newblock \emph{J. Combin. Math. Combin. Comput.}, 40:41--63, 2002.

\bibitem{lubotzky1988}
A.~Lubotzky, R.~Phillips, and P.~Sarnak.
\newblock Ramanujan graphs.
\newblock \emph{Combinatorica}, 8(3):261--277, 1988.

\bibitem{brouwer1989}
A.~E. Brouwer, A.~M. Cohen, and A.~Neumaier.
\newblock \emph{Distance-Regular Graphs}.
\newblock Springer, 1989.

\bibitem{biggs1993}
N.~Biggs.
\newblock \emph{Algebraic Graph Theory}.
\newblock Cambridge University Press, second edition, 1993.

\end{thebibliography}

\end{document}
